\section{KẾ HOẠCH TUẦN TỚI VÀ KẾT LUẬN}\label{sec:ke_hoach}

\subsection{Kế hoạch tuần 02/11/2026 -- 08/11/2026}

Kiểm thử chức năng đã đạt, nên tuần tới tập trung \textbf{đo các chỉ số hiệu năng hoạt động} của từng thành phần trong mạng. Mỗi dòng kế hoạch ứng với một tiêu chí đánh giá của đề cương (Bảng~\ref{tab:tieu_chi}).

\begin{table}[H]
    \centering
    \small
    \caption{Kế hoạch đo chỉ số tuần tới}
    \label{tab:ke_hoach_tuan}
    \begin{tabularx}{\textwidth}{|c|P{3.0cm}|X|P{3.6cm}|}
        \hline
        \textbf{\#} & \textbf{Chỉ số} & \textbf{Cách đo} & \textbf{Kết quả cần đạt} \\
        \hline
        1 & Thời gian khôi phục & SD-WAN: cắt một màu WAN, tắt vSmart, đo lặp 3--5 lần mỗi site, so BFD mặc định với chu kỳ 300~ms; firewall HA đo lặp 3 lần mỗi kịch bản; SDN thử fast-failover Distribution -- Core & Bảng RTO nhỏ nhất/trung bình/lớn nhất và PRR theo thành phần \\
        \hline
        2 & Hiệu năng liên VLAN và liên cơ sở & iperf3 TCP/UDP giữa các VLAN Site 100 và giữa các site; delay, jitter, mất gói; nâng vEdge 2 các site lên 2 vCPU rồi đo lại & Throughput, delay, jitter từng tuyến \\
        \hline
        3 & Tải trên Core và Distribution & CPU, bộ nhớ, lưu lượng interface của Core-SW1/2 và OVS (thống kê cổng qua Ryu) khi chạy tải iperf3 & Biểu đồ tải theo mức lưu lượng \\
        \hline
        4 & Thời gian mở rộng & Bấm giờ thêm một VLAN mới (Ryu, Core, DHCP) và thêm một site vEdge qua vManage & Số phút từ lúc bắt đầu đến khi máy trạm dùng được \\
        \hline
        5 & Quản lý chính sách tập trung & Đổi một chính sách trên vSmart, một luồng trên Ryu hoặc qua lệnh \texttt{/change}; đo thời gian áp xuống toàn mạng & Thời gian lan truyền, số thiết bị chịu ảnh hưởng \\
        \hline
        6 & Độ sẵn sàng & Theo dõi nhiều ngày bằng công cụ web tại Site 900 & Availability, MTTR, MTBF \\
        \hline
    \end{tabularx}
\end{table}

Song song, nhóm đưa số liệu vào chương kết quả của báo cáo cuối kỳ.

\textbf{Lưu ý:} phép đo SDN và SD-WAN không chạy đồng thời để tránh nhiễu số liệu; mỗi kịch bản sự cố đều chụp trạng thái trước và sau làm bằng chứng.

\subsection{Kết luận tuần}

Trong tuần \tuanbaocao, nhóm đã có bộ số liệu SD-WAN đầu tiên (failover link WAN, so sánh hiệu năng năm site), kiểm thử chức năng đạt cho toàn hệ thống, và thiết kế lại cặp firewall HA Site 100: sau tinh chỉnh, hỏng một switch WAN không còn làm Site 100 mất liên lạc với chi nhánh, mọi kịch bản sự cố đều tự phục hồi trong 11--21~s. Kết quả nổi bật khác là Campus NetOps Bot: chạy trong topology, quản trị in-band cả Campus chính lẫn ba chi nhánh qua overlay SD-WAN, chẩn đoán theo lớp, thay đổi cấu hình có duyệt và tự rollback, đạt selftest 23/23. Đối chiếu với đề cương, năm nội dung đã được hiện thực đủ; tuần tới nhóm đo các tiêu chí còn thiếu (thời gian mở rộng, tải Core/Distribution, throughput, quản lý chính sách) trước hạn kiểm thử 10/11.
